\subsection{Integration with respect to the product of two measures}

Throughout this No., \(\lambda\) and \(\mu\) denote measures on S and T, respectively, and \(\nu\) denotes the product measure \(\lambda \otimes \mu\) on \(S \times T\) . In addition, if \(f\) is a positive function on \(S \times T\) , for every \(s \in S\) we denote by \(f_s\) the function \(t \mapsto f(s,t)\) on T, and by \(I_f\) the function \(s \mapsto \mu^\bullet(f_s)\) on S.

Lemma 2. --- Let \(f\) be a \(\nu\) -measurable positive function on \(S \times T\) ; for every compact subset \(L\) of \(T\) , let \(I_f^L\) be the function \(s \mapsto \mu^\bullet(f_s \varphi_L)\) on \(S\) . Then the function \(I_f^L\) is \(\lambda\) -measurable, and

\[
\mathrm{I} _ {f} = \sup _ {\mathrm{L}} \mathrm{I} _ {f} ^ {\mathrm{L}}\tag{5}
\]

\[
\nu^ {\bullet} (f) = \sup _ {L} \lambda^ {\bullet} (I _ {f} ^ {L}),\tag{6}
\]

where L runs over the set of compact subsets of T.

We first note that the inclusion \(L \subset L'\) implies \(I_{f}^{L} \leqslant I_{f}^{L'}\) ; on the other hand, \(\mathrm{I}_{f}^{\mathrm{L}}(s) = \mu_{\mathrm{L}}^{\bullet}((f_{s})_{\mathrm{L}})\) for all \(s \in S\) . Formula (5) is therefore an immediate consequence of the definition of the encumbrance \(\mu^{\bullet}\) given in §1, No.~2. If K is a compact subset of S, and L a compact subset of T, then \(\nu_{\mathrm{K} \times \mathrm{L}} = \lambda_{\mathrm{K}} \otimes \mu_{\mathrm{L}}\) by construction, and Prop. 7 of Ch. V, §8, No.~3 implies the relation

\[
\nu^ {\bullet} (f \varphi_ {\mathrm{K} \times \mathrm{L}}) = \lambda_ {\mathrm{K}} ^ {\bullet} \big ((\mathrm{I} _ {f} ^ {\mathrm{L}}) _ {\mathrm{K}} \big).\tag{7}
\]

Moreover, every compact subset of \(S \times T\) is contained in a compact set of the form \(K \times L\) ; passing to the upper envelope over all K and L in the preceding formula, we therefore obtain

\[
\nu^ {\bullet} (f) = \sup _ {L} \sup _ {K} \lambda_ {K} ^ {\bullet} \big ((I _ {f} ^ {L}) _ {K} \big) = \sup _ {L} \lambda^ {\bullet} (I _ {f} ^ {L}),\tag{8}
\]

namely (6).

Finally, Prop. 7 of Ch. V, §8, No.~3 implies that the restriction of \(I_{f}^{L}\) to every compact subset K of S is \(\lambda_{K}\) -measurable; this is equivalent to saying that \(I_{f}^{L}\) is \(\lambda\) -measurable.

PROPOSITION 11. --- Let \(f\) be a lower semi-continuous function \(\geqslant 0\) defined on \(X = S \times T\) .

\begin{enumerate}
\def\labelenumi{\alph{enumi})}
\item
  The function \(f_{s}: t \mapsto f(s, t)\) is lower semi-continuous on \(\mathbf{T}\) for every \(s \in S\) .
\item
  The function \(I_{f}: s \mapsto \int^{\bullet} f(s, t) d\mu(t)\) is lower semi-continuous on S, and
\end{enumerate}

\[
\iint_ {\mathrm{X}} ^ {\bullet} f (s, t) d \nu (s, t) = \int_ {\mathrm{S}} ^ {\bullet} d \lambda (s) \int_ {\mathrm{T}} ^ {\bullet} f (s, t) d \mu (t).\tag{9}
\]

The property \(a\) ) is obvious, since the mapping \(t \mapsto f(s,t)\) of T into \(\overline{\mathbf{R}}\) is the composition of \(f\) with the continuous mapping \(t \mapsto (s,t)\) of T into X. To establish \(b\) ), we will make use of a lemma:

Lemma 3. --- Let X be a topological space (Hausdorff or not), f a lower semi-continuous function \(\geqslant 0\) defined on X; then f is the limit of an increasing sequence \((f_{n})_{n\in\mathbf{N}}\) of lower semi-continuous functions on X, such that each function \(f_{n}\) is a linear combination, with positive coefficients, of characteristic functions of open sets.

Given two integers \(k \geqslant 1\) and \(n \geqslant 1\) , let us denote by \(J_{kn}\) the characteristic function of the interval \(]k/2^n, +\infty]\) of \(\overline{\mathbf{R}}\) . For every \(x \in \overline{\mathbf{R}}_+\) , set \(u_n(x) = 2^{-n} \sum_{k=1}^{n \cdot 2^n} J_{kn}(x)\) ; it is immediate that the sequence \((u_n(x))_{n \geqslant 1}\) is increasing and admits \(x\) as limit. The sequence of functions \(f_n = u_n \circ f\) is therefore increasing and converges to \(f\) , and one has \(f_{n} = 2^{-n}\sum_{k = 1}^{n\cdot 2^{n}}\varphi_{\mathrm{U}(k,n)}\) ,

where \(\mathrm{U}(k,n)\) is the open set \(\overline{f} (]k / 2^n, + \infty ])\) of X.

Let us pass to the proof of \(b\) ). The function \(I_f\) being the upper envelope of the increasing directed family of functions \(I_f^L\) , where \(L\) runs over the set of compact subsets of \(T\) (Lemma 2), it will suffice to show that the functions \(I_f^L\) are lower semi-continuous; the formula (9) may then be deduced from (6) by passing to the upper envelope over \(L\) (§1, No.~6, Prop. 5).

Thus let \(\mathcal{H}\) be the set of positive lower semi-continuous functions \(f\) on \(\mathrm{S} \times \mathrm{T}\) such that \(\mathrm{I}_f^{\mathrm{L}}\) is lower semi-continuous for every compact subset \(\mathrm{L}\) of \(\mathrm{T}\) . By Prop. 5 of §1, No.~6, the supremum of every increasing directed set of elements of \(\mathcal{H}\) belongs to \(\mathcal{H}\) . By Lemma 3, it will therefore suffice to prove that the characteristic function of an open set \(\mathrm{W}\) of \(\mathrm{S} \times \mathrm{T}\) belongs to \(\mathcal{H}\) . Moreover, by the definition of the product topology on \(\mathrm{S} \times \mathrm{T}\) , the open set \(\mathrm{W}\) is the union of an increasing directed family \((\mathrm{W}_{\alpha})_{\alpha \in \mathrm{A}}\) of open sets of the form

\[
\mathrm{W} = \bigcup_ {1 \leqslant i \leqslant n} \left(\mathrm{U} _ {i} \times \mathrm{V} _ {i}\right),
\]

where the \(U_{i}\) are open in S and the \(V_{i}\) are open in T; by the remarks made above, it will suffice to show that the characteristic function of such an open set belongs to H. Let then \(s \in S\) , and let U be the intersection of the family (possibly empty) formed by the open sets \(U_{i}\) containing s; one sees immediately that \(\varphi_{\mathrm{W}}(s,t) \leqslant \varphi_{\mathrm{W}}(s',t)\) for all \(s' \in U\) and \(t \in T\) , whence, by integration, \(\mathrm{I}_{\varphi_{\mathrm{W}}}^{\mathrm{L}}(s) \leqslant \mathrm{I}_{\varphi_{\mathrm{W}}}^{\mathrm{L}}(s')\) for all \(s' \in U\) . Consequently \(I_{\varphi_{w}}^{L}\) is lower semi-continuous, and the proposition is established.

COROLLARY 1. --- Let \(f\) be a positive numerical function defined on \(\mathbf{X} = \mathbf{S} \times \mathbf{T}\) ; then

\[
\iint_ {\mathrm{X}} ^ {*} f (s, t) d \nu (s, t) \geqslant \int_ {\mathrm{S}} ^ {*} d \lambda (s) \int_ {\mathrm{T}} ^ {*} f (s, t) d \mu (t).\tag{10}
\]

For, let \(g\) be a lower semi-continuous function on \(\mathbf{X}\) such that \(g \geqslant f\) ; by Prop. 11,

\[
\begin{array}{r l} \iint_ {\mathrm{X}} ^ {*} g (s, t) d \nu (s, t) & = \iint^ {\bullet} g (s, t) d \nu (s, t) = \int^ {\bullet} d \lambda (s) \int^ {\bullet} g (s, t) d \mu (t) \\ & = \int^ {*} d \lambda (s) \int^ {*} g (s, t) d \mu (t) \geqslant \int^ {*} d \lambda (s) \int^ {*} f (s, t) d \mu (t). \end{array}
\]

The inequality (10) is obtained by passing to the lower envelope over \(g\) .

COROLLARY 2. --- Let f be a numerical function defined on S × T and ν-negligible. Then the function \(f_{s}: t \mapsto f(s, t)\) is μ-negligible for λ-almost every \(s \in S\) .

PROPOSITION 12. --- Let f be a ν-measurable positive function defined on X = S × T. Assume that f is ν-moderated (resp. that μ is moderated). Then:

\begin{enumerate}
\def\labelenumi{\alph{enumi})}
\item
  The set N of \(s \in S\) such that the function \(f_{s}: t \mapsto f(s, t)\) is not \(\mu\) -measurable is negligible (resp. locally negligible) for \(\lambda\) .
\item
  The mapping \(s \mapsto \int^{\bullet} f(s, t) \, d\mu(t)\) is \(\lambda\) -measurable, and
\end{enumerate}

\[
\iint_ {\mathrm{X}} ^ {\bullet} f (s, t) d \nu (s, t) = \int_ {\mathrm{S}} ^ {\bullet} d \lambda (s) \int_ {\mathrm{T}} ^ {\bullet} f (s, t) d \mu (t).\tag{11}
\]

We begin by establishing \(b)\) when \(f\) is \(\nu\) -moderated. By Lemma 2, this part of the statement is valid when there exists a compact subset \(L\) of \(T\) such that \(f\) is zero outside \(S \times L\) ; for, in this case \(I_f = I_f^{L'}\) for every compact subset \(L'\) of \(T\) containing \(L\) , and formula (11) reduces to (6). In particular, \(b)\) is established for a function \(f\) that is zero outside a compact subset of \(S \times T\) . On the other hand, Cor. 1 of Prop. 11 implies that \(b)\) is true when \(f\) is \(\nu\) -negligible. Since every \(\nu\) -moderated function is the sum of a \(\nu\) -negligible function and a sequence of functions with compact support (§1, No.~9, Cor. 3 of Prop. 14), the assertion \(b)\) is true when \(f\) is \(\nu\) -moderated.

Similarly, the assertion \(b)\) is obvious when \(\mu\) is carried by a compact subset \(L\) of \(T\) (Lemma 2). Suppose that \(\mu\) is moderated; then there exists a sequence \((\mu_n)_{n \in \mathbb{N}}\) of measures on \(T\) with compact support, such that \(\mu = \sum_{n} \mu_n\) (§1, No.~9, Cor. 5 of Prop. 14), whence \(\nu = \sum_{n} \lambda \otimes \mu_n\) (No.~5, Cor. 3 of Prop. 10). The assertion \(b)\) , being valid for each of the measures \(\nu_n = \lambda \otimes \mu_n\) , is also valid for \(\nu = \sum_{n} \nu_n\) .

Let us prove \(a\) ); denote by N the set of \(s \in S\) such that \(f_s\) is not \(\mu\) -measurable; for every compact subset L of T, similarly denote by \(\mathbf{N}_{\mathrm{L}}\) the set of \(s \in S\) such that \(f_s \varphi_{\mathrm{L}}\) is not \(\mu\) -measurable. If K and L are compact sets in S and T respectively, then \(f_{\mathrm{K} \times \mathrm{L}}\) is measurable with respect to the measure \(\nu_{\mathrm{K} \times \mathrm{L}} = \lambda_{\mathrm{K}} \otimes \mu_{\mathrm{L}}\) , and Prop. 2 of Ch. V, §8, No.~2 shows that the set \(\mathbf{N}_{\mathrm{L}}\) is locally negligible for \(\lambda_{\mathrm{K}}\) ; since K is arbitrary, \(\mathbf{N}_{\mathrm{L}}\) is locally \(\lambda\) -negligible.

Suppose that f is zero outside a compact set of the form \(K \times L\) ; then \(N = N_{L}\) , and N is contained in K; it follows that N is \(\lambda\) -negligible. Similarly, if f is \(\nu\) -negligible, Cor. 2 of Prop. 11 implies that N is \(\lambda\) -negligible. The case that f is \(\nu\) -moderated can then be treated as above, on combining the preceding two cases.

Suppose that \(\mu\) is carried by a compact subset L of T; then \(N = N_{L}\) again, therefore N is locally \(\lambda\) -negligible. Since every moderated measure is the sum of a sequence of measures with compact support ( \(\S1\) , No.~9, Cor. 5 of Prop. 14), this result extends at once to the case that \(\mu\) is moderated, on using Prop. 8 of \(\S1\) , No.~7.

Remark. --- Let \((\mathrm{K}_{\alpha})_{\alpha\in\mathrm{A}}\) be a crushing of S for \(\lambda\) and let \(M = S - \bigcup_{\alpha\in A} K_{\alpha}\) ; define in an analogous way \((\mathrm{L}_{\beta})_{\beta\in\mathrm{B}}\) and N for the measure \(\mu\) on T. We denote by \(S'\) the locally compact space that is the sum of the subspaces \(K_{\alpha}\) of S and the discrete space M; the space \(T'\) is defined analogously, and we set \(X' = S' \times T'\) . The locally compact space \(X'\) is the sum of the family \((\mathrm{K}_{\alpha} \times \mathrm{L}_{\beta})_{(\alpha,\beta)\in A\times B}\) of compact subspaces of X and the subspace \(P = (M \times T) \cup (S \times N)\) that is a locally \(\nu\) -negligible subset of X (one observes that P is in general not a discrete space). We saw in the Scholium of §1, No.~8 that there exists a measure \(\lambda'\) on \(S'\) such that the measurable functions, the essential upper integral of positive functions, the essentially integrable functions and their integrals, are the same for \(\lambda\) and \(\lambda'\) . We associate the measure \(\mu'\) on \(T'\) with \(\mu\) , and the measure \(\nu'\) on \(X'\) with \(\nu\) , in conformity with the cited Scholium; one sees immediately that \(\nu'\bullet(f \otimes g) = \lambda'\bullet(f)\mu'\bullet(g)\) for \(f \in \mathcal{F}_{+}(S)\) and \(g \in \mathcal{F}_{+}(T)\) ; therefore \(\nu' = \lambda' \otimes \mu'\) by Prop. 10 of No.~5. Since the topology of \(X'\) is finer than that of X, every \(\nu\) -moderated function is \(\nu'\) -moderated. This procedure permits extending without new proof the Lebesgue--Fubini theorem (Ch. V, §8, No.~4, Th. 1) to the present situation.

